\subsection*{Motivation}

In Chapter 1 we built logic, sets, and proof technique. We now build the arithmetic playground in which all of analysis, optimization, and ultimately neural network training takes place: the real numbers $\mathbb{R}$. The central question of this chapter is not ``what are numbers?'' in a metaphysical sense but rather: \emph{what structural property of $\mathbb{R}$ makes it the right home for limits?} The answer is \textbf{completeness}, and it is what will let us, many chapters later, prove that gradient descent iterates $\theta_t \in \mathbb{R}^d$ actually converge.

\subsection*{Definitions}

We accept $\mathbb{N} = \{0, 1, 2, \ldots\}$ as constructed in Chapter 1 (von Neumann ordinals). The integers $\mathbb{Z}$ extend $\mathbb{N}$ by formal additive inverses; the rationals $\mathbb{Q}$ extend $\mathbb{Z}$ by formal multiplicative inverses of nonzero elements. We then take $\mathbb{R}$ to be a complete ordered field containing $\mathbb{Q}$ (constructed via Dedekind cuts or Cauchy completion of $\mathbb{Q}$). Thus
\[
\mathbb{N} \subset \mathbb{Z} \subset \mathbb{Q} \subset \mathbb{R}.
\]

\begin{definition}[Sequence]
A \emph{sequence} in a set $X$ is a function $a : \mathbb{N} \to X$, written $(a_n)_{n \in \mathbb{N}}$.
\end{definition}

\begin{definition}[Limit, $\varepsilon$--$N$]
A real sequence $(a_n)$ \emph{converges} to $L \in \mathbb{R}$, written $a_n \to L$, iff
\[
\forall \varepsilon > 0 \;\; \exists N \in \mathbb{N} \;\; \forall n \geq N : \;\; |a_n - L| < \varepsilon.
\]
\end{definition}

\begin{definition}[Cauchy]
$(a_n)$ is \emph{Cauchy} iff $\forall \varepsilon > 0\; \exists N\; \forall m, n \geq N : |a_m - a_n| < \varepsilon$.
\end{definition}

\begin{definition}[Bounded, monotone]
$(a_n)$ is \emph{bounded above} iff $\exists M : a_n \leq M$ for all $n$; bounded below symmetrically; \emph{bounded} iff both. It is \emph{monotone increasing} iff $a_n \leq a_{n+1}$ for all $n$.
\end{definition}

\begin{definition}[Supremum, infimum]
Let $S \subseteq \mathbb{R}$ be nonempty and bounded above. $u \in \mathbb{R}$ is an \emph{upper bound} if $s \leq u$ for all $s \in S$. The \emph{supremum} $\sup S$ is the least upper bound: an upper bound such that for every $\varepsilon > 0$ there exists $s \in S$ with $s > \sup S - \varepsilon$. The \emph{infimum} $\inf S$ is defined dually.
\end{definition}

\begin{definition}[Completeness axiom of $\mathbb{R}$]
Every nonempty subset of $\mathbb{R}$ that is bounded above has a supremum in $\mathbb{R}$.
\end{definition}

\subsection*{Theorems and proofs}

\begin{theorem}[Irrationality of $\sqrt{2}$]
There is no rational $q$ with $q^2 = 2$.
\end{theorem}

\begin{proof}
Suppose for contradiction (technique from Chapter 1) that $q = p/r$ with $p, r \in \mathbb{Z}$, $r \neq 0$, the fraction in lowest terms (so $\gcd(p, r) = 1$), and $q^2 = 2$. Then $p^2 = 2 r^2$, so $p^2$ is even, hence $p$ is even (since the square of an odd integer is odd). Write $p = 2k$. Substituting, $4 k^2 = 2 r^2$, i.e. $r^2 = 2 k^2$, so $r^2$ is even and therefore $r$ is even. But then $2 \mid \gcd(p, r)$, contradicting $\gcd(p, r) = 1$.
\end{proof}

\begin{remark}
This is the first concrete witness that $\mathbb{Q}$ is \emph{incomplete}: the bounded-above set $S = \{q \in \mathbb{Q} : q^2 < 2\}$ has no supremum \emph{in $\mathbb{Q}$}. The completeness axiom is exactly the patch that distinguishes $\mathbb{R}$ from $\mathbb{Q}$.
\end{remark}

\begin{theorem}[Bounded monotone convergence]
Let $(a_n)$ be monotone increasing and bounded above in $\mathbb{R}$. Then $(a_n)$ converges, and $\lim a_n = \sup_n a_n$.
\end{theorem}

\begin{proof}
Let $S = \{a_n : n \in \mathbb{N}\}$. $S$ is nonempty (it contains $a_0$) and bounded above by hypothesis. By the completeness axiom, $L := \sup S$ exists in $\mathbb{R}$. Fix $\varepsilon > 0$. By the supremum's least-upper-bound property, $L - \varepsilon$ is \emph{not} an upper bound, so there exists $N \in \mathbb{N}$ with $a_N > L - \varepsilon$. For any $n \geq N$, monotonicity gives $a_n \geq a_N > L - \varepsilon$, while $a_n \leq L$ since $L$ is an upper bound. Hence
\[
L - \varepsilon < a_n \leq L \quad \Longrightarrow \quad |a_n - L| < \varepsilon,
\]
which is exactly the $\varepsilon$--$N$ definition of $\lim a_n = L$.
\end{proof}

\begin{lemma}[Bolzano--Weierstrass for $\mathbb{R}$]
Every bounded real sequence has a convergent subsequence.
\end{lemma}

\begin{proof}
Let $(a_n) \subseteq [c_0, d_0]$ with $c_0 < d_0$. Bisect at midpoint $m_0 = (c_0 + d_0)/2$; at least one of $[c_0, m_0]$, $[m_0, d_0]$ contains $a_n$ for infinitely many $n$ (else only finitely many terms total, contradicting $|\mathbb{N}| = \infty$). Call this half $[c_1, d_1]$ and pick an index $n_1$ with $a_{n_1} \in [c_1, d_1]$. Iterate: at stage $k$ we have nested $[c_k, d_k]$ of length $(d_0 - c_0) / 2^k$ each containing infinitely many terms, and pick $n_k > n_{k-1}$ with $a_{n_k} \in [c_k, d_k]$. The sequence $(c_k)$ is monotone increasing and bounded above by $d_0$, so by Theorem~2.8 it converges to some $L$; similarly $(d_k) \to L$ since $d_k - c_k = (d_0 - c_0)/2^k \to 0$. Since $c_k \leq a_{n_k} \leq d_k$, the squeeze gives $a_{n_k} \to L$.
\end{proof}

\begin{theorem}[Cauchy completeness of $\mathbb{R}$]
Every Cauchy sequence in $\mathbb{R}$ converges.
\end{theorem}

\begin{proof}
Let $(a_n)$ be Cauchy.

\emph{Step 1 (bounded).} Pick $\varepsilon = 1$ and $N$ from the Cauchy condition; for $n \geq N$, $|a_n| \leq |a_N| + 1$, so $(a_n)$ is bounded by $M = \max(|a_0|, \ldots, |a_{N-1}|, |a_N| + 1)$.

\emph{Step 2 (convergent subsequence).} By the Bolzano--Weierstrass lemma, some subsequence $a_{n_k} \to L$.

\emph{Step 3 (full sequence).} Fix $\varepsilon > 0$. Choose $N_1$ with $|a_m - a_n| < \varepsilon/2$ for all $m, n \geq N_1$ (Cauchy property), and choose $K$ with $n_K \geq N_1$ and $|a_{n_K} - L| < \varepsilon/2$ (subsequence convergence). Then for any $n \geq N_1$,
\[
|a_n - L| \leq |a_n - a_{n_K}| + |a_{n_K} - L| < \varepsilon/2 + \varepsilon/2 = \varepsilon. \qedhere
\]
\end{proof}

\subsection*{Code sketch}

Numerically the partial sums $S_n = \sum_{k=1}^n 1/k^2$ approach $\pi^2/6$. We print $|S_n - \pi^2/6|$ and locate the smallest $N$ realizing each tolerance $\varepsilon \in \{10^{-2}, 10^{-4}, 10^{-6}\}$, contrast a Cauchy sequence ($1/n$) with a non-Cauchy one ($(-1)^n$), run a Newton iteration to $\sqrt{2}$, and watch a bounded monotone telescoping series converge to $1$.

\subsection*{Connection to LLMs}

Training a transformer is a sequence $(\theta_t)_{t \in \mathbb{N}}$ in $\mathbb{R}^d$ produced by an optimizer (SGD, Adam). Statements like ``the loss converges'' or ``the iterates converge to a stationary point'' are \emph{limit statements in the sense of Definition 2.2}. Convergence proofs for SGD (Chapter 13) and Adam (Chapter 14) reduce, ultimately, to bounded-monotone arguments on auxiliary scalar sequences (loss values, squared gradient norms) --- and those arguments would simply be false over $\mathbb{Q}$. Completeness is not philosophical decoration; it is the load-bearing axiom under every convergence theorem in deep learning.
